% =====================================================================
%  pds-beamer.tex — Estilo común de las transparencias.
%
%  Uso (desde teoría/Tx/transparencias/):
%     \documentclass[aspectratio=169,11pt]{beamer}
%     \input{../../comun/pds-beamer}
%     \title{...}  \subtitle{Tema x · ...}
%     \begin{document}
%       \portada
%       ...
%     \end{document}
%  Compilar con:  latexmk -xelatex  (o «make» desde teoría/)
% =====================================================================

\usetheme{Boadilla}
\usecolortheme{orchid}
\usefonttheme{professionalfonts}
\setbeamercovered{transparent}
\setbeamersize{text margin left=0.8cm, text margin right=0.8cm}
\setbeamertemplate{navigation symbols}{}
\setbeamertemplate{itemize item}[square]
\setbeamertemplate{itemize subitem}[circle]
\setbeamertemplate{itemize subsubitem}[triangle]
\setbeamertemplate{enumerate item}[square]
\setbeamertemplate{enumerate subitem}[circle]
\setbeamertemplate{caption}{\raggedright\insertcaption\par}

\input{../../comun/pds-comun.tex}

% Tipografía: Fira Sans + Fira Math (legibles en proyector)
\usepackage{fontspec}
\usepackage{unicode-math}
% Se cargan por nombre de fichero para que funcionen con cualquier
% instalación de TeX Live (o MiKTeX) sin instalar fuentes en el sistema.
\setmainfont{FiraSans}[Extension=.otf, UprightFont=*-Regular, BoldFont=*-SemiBold,
  ItalicFont=*-Italic, BoldItalicFont=*-SemiBoldItalic]
\setsansfont{FiraSans}[Extension=.otf, UprightFont=*-Regular, BoldFont=*-SemiBold,
  ItalicFont=*-Italic, BoldItalicFont=*-SemiBoldItalic]
\setmonofont{FiraMono}[Extension=.otf, UprightFont=*-Regular, BoldFont=*-Medium, AutoFakeSlant=0.15, Scale=0.92]
\setmathfont{FiraMath-Regular.otf}
\setmathfont{latinmodern-math.otf}[range={"229B,"22EE,"22F1}]

% Colores del tema ajustados a la paleta común (más contraste)
\setbeamercolor{structure}{fg=pdsazul}
\setbeamercolor{alerted text}{fg=pdsrojo}
\setbeamercolor{example text}{fg=pdsverde}
\setbeamercolor{block title}{bg=pdsazul, fg=white}
\setbeamercolor{block body}{bg=pdsfondo}
\setbeamercolor{block title alerted}{bg=pdsrojo, fg=white}
\setbeamercolor{block body alerted}{bg=pdsrojo!6}
\setbeamercolor{block title example}{bg=pdsverde, fg=white}
\setbeamercolor{block body example}{bg=pdsverde!6}
\setbeamercolor{frametitle}{fg=pdsazul, bg=white}
\setbeamerfont{frametitle}{size=\large, series=\bfseries}
\setbeamerfont{block title}{size=\normalsize, series=\bfseries}
\setbeamertemplate{blocks}[rounded][shadow=false]

% Logo en la esquina superior derecha de cada transparencia
\pgfdeclareimage[height=0.8cm]{logo}{../../comun/ucm.png}
\addtobeamertemplate{frametitle}{}{%
  \begin{tikzpicture}[remember picture, overlay]
    \node[anchor=north east, inner sep=0pt] at ([xshift=-0.25cm, yshift=-0.2cm]current page.north east) {\pgfuseimage{logo}};
  \end{tikzpicture}}

% Datos de portada
\author[\autorescorto]{\autores}
\institute[UCM]{\footnotesize\departamento\\\universidad}
\date[\curso]{\footnotesize\grado\\\curso}

\newcommand{\portada}{%
  \begin{frame}[plain]
    \begin{center}
      \includegraphics[height=1.9cm]{../../comun/ucm.png}
    \end{center}
    \vspace{-0.4cm}
    \titlepage
  \end{frame}}

% Índice al empezar cada sección
\AtBeginSection[]{%
  \begin{frame}{Contenido}
    \tableofcontents[currentsection, hideothersubsections]
  \end{frame}}

\renewcommand{\pdscodesize}{\scriptsize}

% Utilidades
% \herramienta{T1-convolucion-continua}: enlace a la página publicada de
% herramientas/html (se abre en el navegador, sin instalar nada).
% \herramientaservidor{server/....py}: app que necesita Python (bokeh serve).
\newcommand{\urlherramientas}{https://ucm-237.github.io/PdS/}
\newcommand{\herramienta}[1]{%
  \par\vfill\hfill{\scriptsize\color{pdsgris}\faIcon\ \href{\urlherramientas#1.html}{Herramienta interactiva: \texttt{#1}}}}
\newcommand{\herramientaservidor}[1]{%
  \par\vfill\hfill{\scriptsize\color{pdsgris}\faIcon\ Herramienta interactiva (con Python): \texttt{herramientas/#1}}}
\newcommand{\faIcon}{\textbf{▶}}
\newcommand{\gris}[1]{{\color{pdsgris}#1}}
\newcommand{\nota}[1]{\par\vspace{2pt}{\footnotesize\color{pdsgris}\textit{#1}}\par}
% Nota sobre una errata de versiones anteriores (solo con \ERRATAS)
\newcommand{\errata}[1]{\iferratas\nota{Errata de versiones anteriores: #1}\fi}
\newcommand{\apuntes}[1]{\nota{Más detalle en los apuntes, #1.}}

% Bloques con nombre en castellano
\newenvironment{definicion}[1][Definición]{\begin{block}{#1}}{\end{block}}
\newenvironment{importante}[1][Importante]{\begin{alertblock}{#1}}{\end{alertblock}}
\newenvironment{ejemplo}[1][Ejemplo]{\begin{exampleblock}{#1}}{\end{exampleblock}}
